\documentclass[11pt]{article}

\usepackage[margin=1in]{geometry}
\usepackage{amsmath,amssymb,amsthm}
\usepackage{graphicx}
\usepackage[hidelinks]{hyperref}

\newtheorem{theorem}{Theorem}

\newcommand{\en}{e_n}
\newcommand{\lp}{\ell_p}
\newcommand{\ellq}{\ell_q}
\newcommand{\Lp}{L_p}

\begin{document}

\begin{center}
{\Large Literature status: the \(p,q>2\) Lipschitz extension upper bound}

\medskip
{\small Run \texttt{fa\_banach\_001}; source arXiv:1506.04398}
\end{center}

\section*{Source question}

Naor and Rabani's paper \emph{On Lipschitz extension from finite subsets}
defines \(\en(X,Z)\) as the supremal Lipschitz extension loss for
Banach-space-valued functions from subsets of \(X\) of cardinality at most
\(n\).  In Section 1.4 they survey open questions about
\(\en(\ell_p,\ell_q)\).  For \(p,q\in(2,\infty)\), the only upper bound that
they record is the general Lee--Naor finite-metric bound
\[
  \en(\ell_p,\ell_q)\lesssim \frac{\log n}{\log\log n},
\]
and they explicitly say that no asymptotically smaller upper bound was known.

\begin{figure}[h]
\centering
\includegraphics[width=\textwidth]{figures/open_problem_crop.png}
\caption{arXiv:1506.04398, page 8: the source passage recording the
\(p,q>2\) upper-bound gap for \(\en(\ell_p,\ell_q)\).}
\end{figure}

\section*{Supporting literature status}

Naor and Ren's later paper \emph{Euclidean embedding, randomized clustering,
and Lipschitz extension for finite and doubling subsets of \(L_p\) when
\(p>2\)} proves a stronger Banach-target extension theorem for domains
\(L_p\).  They also explicitly state that the theorem answers a natural
question left open by \cite{LN05} and stated in \cite{NR17}.

\begin{figure}[h]
\centering
\includegraphics[width=\textwidth]{figures/supporting_answer_crop.png}
\caption{arXiv:2502.10543, page 7: Theorem 1.5 and the explicit answer
statement.}
\end{figure}

\begin{figure}[h]
\centering

\includegraphics[width=\textwidth]{figures/supporting_reference_crop.png}
\caption{arXiv:2502.10543, page 59: bibliography entry [82], identifying
\cite{NR17} as Naor--Rabani's \emph{On Lipschitz extension from finite
subsets}.}
\end{figure}

\section*{Proof intuition}

The source passage asks whether finite subsets of \(\ell_p\), with
\(p>2\), have any extension behavior better than arbitrary finite metric
spaces when the target is another \(\ell_q\), \(q>2\).  Naor--Ren prove
something stronger at the domain level: every \(n\)-point subset of \(L_p\)
has Banach-valued Lipschitz extension loss \(O_p(\sqrt{\log n})\).  Since
\(\ell_p\) embeds isometrically into an \(L_p\) space, an extension theorem
on \(L_p\) immediately restricts to \(\ell_p\).  Then choosing the target
Banach space to be \(\ell_q\) gives exactly the source paper's
\(p,q>2\) case.

\section*{The literature answer}

\begin{theorem}
For every fixed \(p,q\in(2,\infty)\),
\[
  \en(\ell_p,\ell_q)\lesssim_p \sqrt{\log n}.
\]
In particular,
\[
  \en(\ell_p,\ell_q)
  = o\!\left(\frac{\log n}{\log\log n}\right)
  \qquad (n\to\infty).
\]
Thus the \(p,q>2\) upper-bound gap recorded in \cite{NR17} has been answered
by \cite{NRen}.
\end{theorem}

\begin{proof}
Naor--Ren's Theorem 1.5 \cite{NRen} states that if \(2<p<\infty\), then
\[
  \en(L_p)\lesssim p^2\sqrt{\log n}.
\]
Here \(\en(L_p)\) is the supremum over all subsets of \(L_p\) with
cardinality at most \(n\) and over all Banach targets.

Embed \(\ell_p\) isometrically into an \(L_p\) space; for instance, use
disjoint measurable sets and represent a sequence by the corresponding
step function after the standard normalization.  Let
\(S\subset \ell_p\) with \(|S|\le n\), and let
\(f:S\to \ell_q\) be Lipschitz.  Viewing \(S\) as a subset of \(L_p\),
Theorem 1.5 of \cite{NRen} extends \(f\) to a map
\(F:L_p\to\ell_q\) with
\[
  \|F\|_{\operatorname{Lip}}\lesssim p^2\sqrt{\log n}\,
  \|f\|_{\operatorname{Lip}}.
\]
Restricting \(F\) back to the embedded copy of \(\ell_p\) gives an extension
of \(f\) to \(\ell_p\) with the same Lipschitz bound.  Taking the supremum
over \(S\) and \(f\) yields
\[
  \en(\ell_p,\ell_q)\lesssim p^2\sqrt{\log n}.
\]
For fixed \(p\), \(p^2\sqrt{\log n}=o((\log n)/\log\log n)\), so this is
asymptotically smaller than the general \cite{LN05} bound recorded in the
source paper.
\end{proof}

\section*{Verification notes and limitations}

This packet is a literature-status identification, not a new proof from the
run.  The supporting paper explicitly says that Theorem 1.5 answers the
natural question left open by \cite{LN05} and stated in \cite{NR17}; its
bibliography identifies \cite{NR17} as Naor--Rabani's 2017 journal version of
arXiv:1506.04398.

The packet addresses the source paper's \(p,q>2\) upper-bound gap for
\(\en(\ell_p,\ell_q)\).  It does not answer the separate Kalton
Holder-extension questions or the finite-dimensional Ball extension question
that also appear in arXiv:1506.04398.  It also does not assert sharp
asymptotics for \(\en(\ell_p,\ell_q)\); it records that an upper bound
strictly below \(O((\log n)/\log\log n)\) is now known.

Before promotion, the run lightweight indexes were searched for arXiv:1506.04398,
the title phrase, ``Lipschitz extension'', ``finite subsets'', ``ell\_p'',
``ell\_q'', and nearby core terms.  No existing exact packet or attempt was
found.  Broad hits to other Lipschitz/free-space packets were not duplicates,
and human-only archive contents were not opened.

\section*{Human review recommendation}

Treat the \(p,q>2\) upper-bound gap recorded in arXiv:1506.04398 Section 1.4
as already answered by arXiv:2502.10543.  Keep the named Kalton and Ball
questions from the same source paper separate.

\begin{thebibliography}{9}

\bibitem{LN05}
J. R. Lee and A. Naor,
\emph{Extending Lipschitz functions via random metric partitions},
Invent. Math. 160 (2005), 59--95.

\bibitem{NR17}
A. Naor and Y. Rabani,
\emph{On Lipschitz extension from finite subsets},
Israel J. Math. 219 (2017), 115--161; arXiv:1506.04398.

\bibitem{NRen}
A. Naor and K. Ren,
\emph{Euclidean embedding, randomized clustering, and Lipschitz extension for
finite and doubling subsets of \(L_p\) when \(p>2\)},
arXiv:2502.10543.

\end{thebibliography}

\end{document}
